\section{Rigidity of the seven-cubic banks and a one-pole obstruction}
\label{app:seven-bank-augmentation}

We work in characteristic zero.  A seven-cubic bank means seven distinct
sections of $\mathcal O_{\mathbf P^1}(3)$ agreeing with a word on fourteen
distinct points in at least seven positions each.  Equivalence permits
projective changes of the evaluation coordinate, a common nonzero scaling,
and addition of a common cubic section.

\begin{proposition}\label{prop:seven-bank-rigidity}
Every seven-cubic bank is equivalent to a Galois conjugate of one of the
two banks of incidence types 2 and 7 in the accompanying incidence
certificate.  In particular, type 2 has no additional realization
component or nontrivial normalized infinitesimal deformation.
\end{proposition}
\begin{proof}
Pair counting gives at least $63$ common pair agreements and at most
$3\binom72=63$.  Equality forces seven triple columns, seven quadruple
columns, seven agreements per candidate, and three per pair.  Writing
$T$ for the triple columns and $C$ for the complements of the quadruple
columns gives the eight ordered incidence types enumerated in
Appendix~\ref{app:global-seven-cubics}.  The guard-preserving identities
in \texttt{NONFANO\_ORBITS\_01\_OBSTRUCTION.md},
\texttt{ORBIT3\_OBSTRUCTION.md}, and the type-4--6 proof of that
appendix exclude all but types 2 and 7.
The normalized type-7 classification is proved in that appendix.
Here we complete the type-2 classification.

Its ordered blocks are
\[
 C=(137,145,167,236,247,256,345),\qquad
 T=(123,124,156,257,346,357,467).
\]
Normalize the quadruple coordinates to $\infty,0,1,u,v,z,d$.
After subtracting a common cubic, all quadruple-compatible families are
obtained by reducing
$W=aX^3+bX^4+cX^5$ modulo each candidate's finite quadruple locator.
Indeed, set the infinity word to zero, interpolate the other six values
by a quintic, and subtract its quadratic part.  A candidate matching
infinity has a cubic locator and a quadratic remainder; the others
have quartic locators and cubic remainders.  The amplitude vector is
nonzero, since the candidates are distinct.

The seven triple conditions give seven linear equations in $(a,b,c)$.
The $156$ equation has $a$-coefficient $-1$, so eliminate $a$ without
a denominator.  A two-by-two minor of the resulting six-by-two matrix is
\[
 -(d-1)(d-uv+uz-z)(-dvz+dv+uvz-uv+uz-vz).
\]
Thus there are two branches.  On $d=uv-uz+z$, another equation gives
$u=vz/D$, where $D=v+z^2-vz\ne0$.  A remaining minor is
\[
 \frac{vz(v-z)^2(v+z^2-2vz)}D.
\]
It is nonzero: its last factor divided by $D$ is $1-u$.
This branch is impossible.

On the other branch,
\[
 d=\frac{uv(z-1)+z(u-v)}{v(z-1)},\qquad
 u=\frac{vz^2(v-z)}{D_1},\quad
 D_1=v^2z-v^2-vz^3+2vz^2-vz+v-z^2.
\]
The first denominator is a node guard; the second cannot vanish because
its defining equation would give $vz^2(v-z)=0$.
Two necessary remaining equations are
\begin{align*}
 E&=v^2z-2v^2+vz+v-z^2,\\
 F&=v^3z^3-v^3z^2+v^3z-v^3-2v^2z^3
       +2v^2z^2+v^2z-2vz^3+z^5.
\end{align*}
Direct polynomial elimination gives
\[
 \operatorname{Res}_v(E,F)=z^6(z-1)^4f(z),\qquad
 f(z)=z^3-10z^2+3z+1.
\]
Since $z\ne0,1$, we have $f(z)=0$.  Reduction of $E,F$ on this
cubic gives $25v=-4z^2+27z+7$, and hence
\[
 u=\frac{-z^2+13z+3}{20},\qquad
 v=\frac{-4z^2+27z+7}{25},\qquad
 d=\frac{-3z^2+29z+4}{20}.
\]
The amplitude matrix has rank two: its rows $0,1$, columns $0,1$ minor
is $1$ modulo the good prime $z\mapsto4$ over $83$, and the archived
exact bank supplies a nonzero kernel vector over $\mathbf Q[z]/(f)$.
That vector has $a\ne0$, so it is unique after $a=1$.
Each triple contains a pair belonging to no other triple; its two
quadruple roots leave a unique possible third root.  Thus the triple
coordinates are also fixed.  The archived exact bank verifies existence
and all guards at the three conjugates.

For completeness, the seven normalized equations in
$(u,v,z,d,b,c)$ have a six-by-six Jacobian minor equal to $14$ at
$(31,12,4,70,69,25)$ modulo $83$.  The normalized tangent space is
therefore zero in characteristic zero.  The exact resultants,
reductions, rank minor, and Jacobian are independently replayed by
\texttt{orbit2\_rigidity\_audit.py}; no modular nonlinear
inconsistency inference is used.
\end{proof}

\begin{theorem}\label{thm:seven-bank-one-pole-obstruction}
Let a seven-cubic bank be given over a characteristic-zero field.
Let $\psi=[R:S]:\mathbf P^1\to\mathbf P^1$ be a rational map of
degree $e\ge2$, with coprime binary forms $R,S$ of degree $e$.
Assume the fourteen selected fibers are separable and disjoint, and
pull back the word by the section multiplier $S^3$.
There is no binary form $N$ of degree $3e+1$ and linear form $\ell$
whose zero $b$ is outside the selected fibers, with $N(b)\ne0$, such
that $N/\ell$ has at least $7e$ agreements with this pulled-back word.
The pole at infinity is included.
\end{theorem}

This excludes the stated proper one-pole augmentation of a seven-cubic
bank, not arbitrary eight-word constructions or arbitrary covers of
larger-degree banks.

\begin{proof}
We may extend constants algebraically and use
Proposition~\ref{prop:seven-bank-rigidity}.  All transformations in that
proposition preserve the asserted construction and its properness.
We first give the cover-independent norm bound underlying the finite
linear certificates.

\smallskip\noindent\textbf{Pole and norm bound.}
Suppose $f$ selected fibers are fully matched.  Let $L$ interpolate
their base word values, with $\deg L\le f-1$, and let $Q$ be their
monic locator; for $f=0$ use $L=0,Q=1$.  Put
\[
 v=\frac{N/(\ell S^3)-L(\psi)}{Q(\psi)},\qquad m=3-f.
\]
Full-fiber cancellation gives
\[
 (v)_\infty\le m\psi^*(\infty)+[b].                 \tag{W1}
\]
The pair $(\psi,v)$ is birational onto its image, even when $e$ is
composite.  Otherwise both functions descend through a finite map
$\phi$ of degree $t>1$ to the normalization of the image.  The pole
excess of $v$ over $m$ times the pole order of $\psi$ is exactly one
at $b$, and is nonpositive everywhere else by (W1).  Pole orders pull
back by ramification indices.  Thus every point of the fiber
$\phi^{-1}(\phi(b))$ would have positive excess.  This fiber must
consist only of $b$, whose ramification index is then $t$, contradicting
that its excess is one.  The exact excess at $b$ survives subtraction
and division by the base interpolant and locator, because $b$ is
outside the selected fibers.

The monic minimal equation of $v$ has degree $e$.  At finite places its
symmetric coefficients have no poles except possibly a simple pole at
$\psi(b)$: the sum of the additional negative conjugate valuations is
at most one.  At infinity the $j$th symmetric coefficient has pole
order at most $jm$, with an additional at most one if $\psi(b)=\infty$.
Clearing the possible finite pole therefore gives a nonzero primitive
relation
\[
 H(X,V)=\sum_{j=0}^e A_j(X)V^{e-j},\qquad
 \deg A_j\le j(3-f)+1,
 \quad A_0\ne0.                                      \tag{W2}
\]
This argument includes ramification and cancellation at poles.
At a remaining matched fiber the value of $v$ is the fixed scalar
$y_x=(w_x-L(x))/Q(x)$.  Two distinct matched preimages imply
\[
 H(x,y_x)=H_X(x,y_x)=H_V(x,y_x)=0,                     \tag{W3}
\]
since the birational image then has two normalization branches there.
Only the value equation is imposed at a single match.


There is also a uniform singularity budget.  Compactify the total space
of $\mathcal O_{\mathbf P^1}(m)$ to the Hirzebruch surface $\mathbb F_m$.
Write $C_0$ for its negative section and $F$ for a fiber, so
$C_0^2=-m$, $C_0F=1$, and $K=-2C_0-(m+2)F$.
The irreducible closure of the primitive equation (W2) has class
\[
 D=eC_0+(me+c)F,\qquad
 c=\max_j(\deg A_j-mj)\in\{0,1\}.
\]
Indeed these coefficient degrees are precisely the transition rule
for the fiber coordinate of $\mathcal O(m)$; choosing the actual
maximum removes any artificially added fiber at infinity.  Primitivity
removes vertical components, and the nonzero leading coefficient makes
$c\ge0$.  Adjunction gives
\[
 p_a(D)=1+\frac{D(D+K)}2
       =(e-1)\left(\frac{me}{2}+c-1\right)
       \le\frac{me(e-1)}2.                            \tag{W4}
\]
Its normalization is $\mathbf P^1$.  Hence the sum of finite
$\delta$-invariants is at most (W4), with no nonsingularity assumption
at the boundary.  If $k_x$ distinct selected preimages match at $x$,
they give $k_x$ distinct branches at $(x,y_x)$, so that point contributes
at least $\binom{k_x}{2}$.

\smallskip\noindent\textbf{All degrees at least four.}
Select exactly $7e$ matches, and let $T,Q$ count those above the triple
and quadruple base nodes.  Each old cleared section differs from $N$
by a nonzero degree-$(3e+1)$ binary form, since the witness is proper.
Summing its root bound over the seven old candidates gives
\[
 T+Q=7e,\qquad 3T+4Q\le7(3e+1),\qquad T\ge7(e-1).
\]
Convexity on the seven triple fibers and (W4), with $m=3$, imply
\[
 \frac{7(e-1)(e-2)}2
 \le\sum_{x\text{ triple}}\binom{k_x}{2}
 \le\frac{3e(e-1)}2.
\]
This is impossible for every integer $e\ge4$.

\smallskip\noindent\textbf{Quadratic covers.}
Move $b$ to source infinity.  The proper witness becomes a polynomial
of exact degree seven, while the old sections have degree at most six.
Four full fibers are impossible: subtract the cleared cubic interpolant
through their values to obtain a degree-at-most-seven polynomial with
eight roots.  Select exactly fourteen matches.  Full and empty fibers
form disjoint sets $F,E$ of the same size $f\le3$.
Writing $I_i(x)$ for the base agreement indicator of candidate $i$,
pairwise degree-seven root counting implies
\[
 \sum_{x\in F}I_i(x)\le\sum_{x\in E}I_i(x)
 \quad(1\le i\le7).                                  \tag{W5}
\]
Enumerating these finite subsets gives the counts in the table below.
The relation (W2) has coefficient degree caps $(1,4-f,7-2f)$ and is
imposed at $14-2f$ single fibers.
\[
\begin{array}{c|rrrr|rrrr}
 &\multicolumn{4}{c|}{\text{number of patterns}}&
   \multicolumn{4}{c}{\text{matrix columns}}\\
 f&0&1&2&3&0&1&2&3\\\hline
 \text{type 7}&1&7&42&252&15&12&9&6\\
 \text{type 2}&1&1&54&276&15&12&9&6
\end{array}                                             \tag{W6}
\]
All $f\ge1$ matrices have full column rank.  For $f=0$ both
fourteen-by-fifteen matrices have rank fourteen.  For type 7 the
received word has the form $(a+cX^7)/X^2$ in its Paley chart; hence the
kernel is spanned by $X^2V-a-cX^7$, contradicting $A_0\ne0$.
For type 2 the unique kernel relation is quadratic, but its discriminant
has degree eight and is squarefree.  It defines a function field of
genus three, contradicting $k(X,v)=k(U)$.  These exact rank and
squarefreeness certifications are detailed below.

\smallskip\noindent\textbf{Cubic covers: the fiber budget.}
Let $d$ and $f$ denote the numbers of double and full matched fibers.
Formula (W4), with $e=3,m=3$, gives $d+3f\le9$.
It forces $f\le2$ for a twenty-one-match witness: $f=3$ would force
$d=0$ and hence at most twenty matches, and $f\ge4$ is impossible.
If $f=2$, subtract the linear interpolant and divide by the two full
fiber locators.  Applying (W4) with $m=1$ gives $d\le3$.
This argument includes every ramification partition above infinity,
cancellation at the poles, and the extra pole at infinity itself.

Combining these inequalities with the seven pair intersection bounds
leaves just two forms of support: every triple fiber double and every
quadruple fiber single; or one triple fiber full, the other six double,
one quadruple fiber empty, and the other six single.  There are seven
allowed full/empty pairs for type 7 and one for type 2.
This finite assertion can be checked directly from the displayed
incidence blocks: their triple and quadruple incidence matrices are
invertible; two full fibers would force their nonconstant excess
vector to be constant.  The remaining one-full alternatives reduce to
column differences of the two incidence matrices and fail unless the
full triple lies in the empty quadruple.  The complete integer census
is also supplied with the certificates.

For the no-full pattern, (W2)--(W3) give $28$ rows and $26$ columns,
with coefficient caps $(1,4,7,10)$.  For a one-full pattern they give
$24$ rows and $20$ columns, with caps $(1,3,5,7)$.
Every one of these matrices has full column rank, contradicting (W2).
This completes the proof.
\end{proof}

\paragraph{Exact certificate provenance.}
The finite linear assertions above use reductions of the specified
number-field banks, not a search over covers.  For type 7 use the
Paley chart with $\zeta_7\mapsto16$ in $\mathbf F_{29}$; for type 2
use $z\mapsto4$ in $\mathbf F_{83}$ and the independently verified
affine bank.  Every node difference and interpolation denominator is
a unit.  The quadratic certificates contain all $302$ and $332$
matrices and selected minors; independent verifiers rebuild the
incidence masks, enumerate (W5), and evaluate these minors by integer
Bareiss elimination.  The cubic certificates contain eight type-7
matrices and two type-2 matrices; the latter maximal minors are
$38$ and $1$ modulo $83$.  Separate verifiers reconstruct (W2)--(W3).

In the type-2 quadratic no-full case, normalize the free kernel
coordinate to one.  Its discriminant reduces to
\[
 -8X^8+35X^7+36X^6+28X^5-13X^4-21X^3-3X^2+24X+36
 \quad\text{in }\mathbf F_{83}[X].
\]
The supplied Bezout identity with its derivative proves squarefreeness.
The unit rank minor makes the normalized characteristic-zero kernel
integral at this prime by Cramer's rule.  Thus degree eight and
squarefreeness lift, as do all nonzero minors.  Scalar extension and
Galois conjugation preserve the conclusions.

The files are in \texttt{quadratic\_one\_pole\_route} and
\texttt{cubic\_cover\_norm\_gate}, with independent verification
receipts and the source bank arrays.  The orbit-2 elimination and
Jacobian are in \texttt{fano\_seven/orbit2\_rigidity\_audit.json};
its branch formulas are replayed from
\texttt{nonfano2\_projective.json} and the two branch files.
All lifting arguments here are nonvanishing-minor or polynomial-identity
arguments; none infers characteristic-zero inconsistency from a
modular Groebner computation.
